\section{Introdução}
A análise de partidas de futebol tem se beneficiado cada vez mais do uso de técnicas computacionais para extrair informações diretamente de vídeos. A partir das imagens de uma partida, é possível obter dados sobre o posicionamento dos jogadores, movimentações, ocupação do espaço, organização das equipes e ocorrência de diferentes ações de jogo. Essas informações podem ser utilizadas tanto para análises táticas quanto para o desenvolvimento de sistemas automatizados de acompanhamento e reconstrução do estado da partida.\\

Nesse contexto, técnicas de processamento digital de sinais e visão computacional desempenham um papel importante, uma vez que permitem transformar sequências de imagens em informações quantitativas sobre os elementos presentes no campo. Entretanto, a análise de vídeos de futebol apresenta desafios particulares. Os jogadores estão constantemente em movimento, podem sofrer oclusões, possuem aparência semelhante entre si e a câmera pode apresentar diferentes perspectivas e movimentos. Além disso, uma representação baseada apenas nas coordenadas dos jogadores na imagem não corresponde diretamente à sua posição no campo, devido à perspectiva da câmera.\\

Uma abordagem para esse problema consiste em realizar a reconstrução do estado do jogo (Game State Reconstruction — GSR), relacionando os elementos observados no vídeo às suas posições em um sistema de coordenadas associado ao campo. Essa representação permite analisar a distribuição espacial dos jogadores de maneira independente da perspectiva da câmera, tornando possível estudar aspectos como posicionamento e ocupação do campo. Para isso, informações geométricas da imagem, como a relação entre o plano da imagem e o plano do campo, podem ser obtidas por meio de técnicas como a homografia.\\

Neste projeto, é explorado o problema de identificar a posição de jogadores de futebol a partir de uma sequência de vídeo, utilizando como base vídeos panorâmicos de partidas completas e as informações de referência disponibilizadas pelo conjunto de dados SoccerTrack v2. O conjunto de dados foi desenvolvido especificamente para tarefas de reconstrução do estado da partida e reúne vídeos de partidas completas, acompanhados de anotações que incluem coordenadas dos jogadores no campo, identificadores individuais, números das camisas, funções e equipes. Essa estrutura fornece as informações necessárias para investigar uma abordagem de processamento de vídeo voltada à obtenção da posição dos jogadores em coordenadas do campo.\\

A proposta deste trabalho é, portanto, desenvolver uma aplicação capaz de receber um vídeo de uma partida de futebol e processar seus quadros para estimar a posição dos jogadores no campo e associá-los às respectivas equipes. Para isso, são exploradas etapas de processamento e análise das imagens, juntamente com informações geométricas da câmera e do campo. Como resultado, busca-se demonstrar a aplicação de técnicas de processamento digital de sinais e visão computacional a um problema concreto de análise esportiva, produzindo uma representação mais adequada para a análise espacial dos jogadores durante a partida.\\